\documentclass[11pt]{article}
\usepackage{amsmath, amssymb, amsthm}
\usepackage[margin=1in]{geometry}

\newtheorem{theorem}{Theorem}[section]
\newtheorem{proposition}[theorem]{Proposition}
\newtheorem{lemma}[theorem]{Lemma}
\newtheorem{corollary}[theorem]{Corollary}
\newtheorem{definition}[theorem]{Definition}
\newtheorem{remark}[theorem]{Remark}
\newtheorem{question}[theorem]{Question}

\title{Period Polynomials and Special L-Values at Weight 12:\\
A Synthesis of Classical and Weakly Holomorphic Theory\\
{\large Extended Version with BFK Computational Formulas}}
\date{November 28, 2025 (Expanded)}

\begin{document}

\maketitle

\begin{abstract}
We synthesize three perspectives on periods of modular forms at weight 12: Zagier's classical theory of period polynomials for cusp forms, the Bringmann-Fricke-Kent framework for L-values of weakly holomorphic modular forms, and explicit rationality patterns observed in special L-values of the discriminant form $\Delta$. 

This expanded version adds explicit computational formulas using the BFK incomplete gamma framework, providing rapidly convergent expressions for testing rationality patterns in both $\Delta$ (the unique weight 12 cusp form) and its weakly holomorphic Hecke partner $\Delta'$ (with principal part $q^{-1}$).
\end{abstract}

\tableofcontents

\section{Introduction}

The theory of periods of modular forms connects special values of L-functions, period integrals, and the cohomology of modular curves. At weight 12, where $\dim S_{12} = 1$, the unique normalized cusp form
\[
\Delta(z) = q \prod_{n=1}^{\infty} (1-q^n)^{24} = \sum_{n=1}^{\infty} \tau(n) q^n
\]
provides the simplest non-trivial example where these connections can be made completely explicit.

\subsection{Main Objects of Study}

\begin{definition}
The \textbf{classical discriminant} is the unique normalized cusp form of weight 12:
\[
\Delta(z) = q - 24q^2 + 252q^3 - 1472q^4 + \cdots \in S_{12}.
\]
\end{definition}

\begin{definition}
The \textbf{weakly holomorphic partner} $\Delta'$ is the unique weakly holomorphic modular form of weight 12 with principal part $q^{-1}$:
\[
\Delta'(z) = q^{-1} + 0 + 0 + c_2 q^2 + c_3 q^3 + \cdots \in M_{12}^!.
\]
Note the gap: the first non-zero Fourier coefficient after the pole occurs at $q^2$, not $q^1$. This is a defining property of the Duke-Jenkins basis element $f_{12,1}$.
\end{definition}

These forms are related by Hecke duality through the Duke-Jenkins basis structure for weakly holomorphic modular forms.

\section{Zagier's Period Polynomial Theory}

\subsection{Period Polynomials}

\begin{definition}[Zagier]
For a cusp form $f \in S_k$ of weight $k$ on $\mathrm{PSL}_2(\mathbb{Z})$, the \textbf{period polynomial} is
\[
r_f(X) := \int_0^{i\infty} f(z)(z-X)^{k-2} \, dz = -\sum_{n=0}^{k-2} \frac{(k-2)!}{(k-2-n)!} \frac{L(f,n+1)}{(2\pi i)^{n+1}} X^{k-2-n}.
\]
\end{definition}

The period polynomial encodes all critical L-values $L(f,1), L(f,2), \ldots, L(f,k-1)$ in the critical strip.

\subsection{Decomposition into Even and Odd Parts}

Every period polynomial decomposes as $r_f(X) = r_f^+(X) + r_f^-(X)$ where:
\begin{align*}
r_f^+(X) &= \frac{1}{2}(r_f(X) + r_f(-X)) \quad \text{(even part)}, \\
r_f^-(X) &= \frac{1}{2}(r_f(X) - r_f(-X)) \quad \text{(odd part)}.
\end{align*}

\subsection{Basis Polynomials at Weight 12}

\begin{definition}
The \textbf{template polynomials} at weight 12 are:
\begin{align*}
P_0(X) &= X^{10} - 1, \\
P_1(X) &= X^8 - 3X^6 + 3X^4 - X^2, \\
P_2(X) &= 4X^9 - 25X^7 + 42X^5 - 25X^3 + 4X.
\end{align*}
\end{definition}

\subsection{Explicit Decomposition for $\Delta$}

\begin{theorem}[Zagier]
The period polynomial of $\Delta$ decomposes as:
\begin{align*}
r_\Delta^+(X) &= \omega_+ \left( -\frac{192}{691} P_0(X) + \frac{16}{3} P_1(X) \right), \\
r_\Delta^-(X) &= \omega_- \cdot 192 \cdot P_2(X),
\end{align*}
where $\omega_\pm$ are the fundamental periods of $\Delta$:
\begin{align*}
\omega_+ &= \frac{(k-2)!}{(4\pi)^{k-1}} \zeta(k-1) = \frac{10!}{(4\pi)^{11}} \zeta(11), \\
\omega_- &= \frac{\omega_+}{i \langle \Delta, \Delta \rangle}.
\end{align*}
\end{theorem}

\section{Rationality Patterns in Special L-Values}

\subsection{The Lockstep Rationality Phenomenon}

\begin{theorem}[Lockstep Rationality]
After normalizing by $(2\pi i)^{n+1}$ and factorials, the special L-values of $\Delta$ satisfy:
\[
\frac{L(\Delta, n+1)}{(2\pi i)^{n+1} \cdot n!} \in \mathbb{Q} \cdot \omega_+ \quad \text{for even } n,
\]
\[
\frac{L(\Delta, n+1)}{(2\pi i)^{n+1} \cdot n!} \in \mathbb{Q} \cdot \omega_- \quad \text{for odd } n.
\]
\end{theorem}

\subsection{L-Values at Even $s$ (from $P_2$ and $\omega_-$)}

From $r_\Delta^-(X) = 192\omega_- \cdot P_2(X)$ with $P_2(X) = 4X^9 - 25X^7 + 42X^5 - 25X^3 + 4X$:

\begin{proposition}
The special L-values of $\Delta$ at even $s$ are:
\begin{align*}
\frac{L(\Delta, 10)}{(2\pi i)^{10} \cdot 9! \cdot \omega_-} &= 192 \cdot 4 = 768, \\
\frac{L(\Delta, 8)}{(2\pi i)^8 \cdot 7! \cdot \omega_-} &= 192 \cdot (-25) = -4800, \\
\frac{L(\Delta, 6)}{(2\pi i)^6 \cdot 5! \cdot \omega_-} &= 192 \cdot 42 = 8064, \\
\frac{L(\Delta, 4)}{(2\pi i)^4 \cdot 3! \cdot \omega_-} &= 192 \cdot (-25) = -4800, \\
\frac{L(\Delta, 2)}{(2\pi i)^2 \cdot 1! \cdot \omega_-} &= 192 \cdot 4 = 768.
\end{align*}
\end{proposition}

\subsection{L-Values at Odd $s$ (from $P_0, P_1$, and $\omega_+$)}

From $r_\Delta^+(X) = \omega_+ \left( -\frac{192}{691} P_0(X) + \frac{16}{3} P_1(X) \right)$:

\begin{proposition}[Contribution from $P_0$]
The polynomial $P_0(X) = X^{10} - 1$ contributes to $L(\Delta,1)$ and $L(\Delta,11)$:
\begin{align*}
\frac{L(\Delta, 11)}{(2\pi i)^{11} \cdot 10! \cdot \omega_+} &= -\frac{192}{691}, \\
\frac{L(\Delta, 1)}{(2\pi i)^1 \cdot 0! \cdot \omega_+} &= \frac{192}{691}.
\end{align*}
\end{proposition}

\begin{proposition}[Contribution from $P_1$]
The polynomial $P_1(X) = X^8 - 3X^6 + 3X^4 - X^2$ contributes:
\begin{align*}
\frac{L(\Delta, 9)}{(2\pi i)^9 \cdot 8! \cdot \omega_+} &= \frac{16}{3}, \\
\frac{L(\Delta, 7)}{(2\pi i)^7 \cdot 6! \cdot \omega_+} &= -16, \\
\frac{L(\Delta, 5)}{(2\pi i)^5 \cdot 4! \cdot \omega_+} &= 16, \\
\frac{L(\Delta, 3)}{(2\pi i)^3 \cdot 2! \cdot \omega_+} &= -\frac{16}{3}.
\end{align*}
\end{proposition}

\section{BFK Computational Framework for $\Delta$}

\subsection{Rapidly Convergent L-Function Expressions}

The BFK paper provides expressions for L-functions using incomplete gamma functions that converge exponentially fast. For a cusp form $f \in S_k$ with Fourier expansion $f(z) = \sum_{n=1}^{\infty} a_f(n) q^n$, the completed L-function is:

\begin{theorem}[BFK Formula for Cusp Forms]
For any $t_0 > 0$:
\[
L^*_f(s) = \sum_{n=1}^{\infty} a_f(n) \frac{\Gamma(s, 2\pi n t_0)}{(2\pi n)^s} + i^k \sum_{n=1}^{\infty} a_f(n) \frac{\Gamma(k-s, 2\pi n/t_0)}{(2\pi n)^{k-s}},
\]
where $\Gamma(s,x) = \int_x^{\infty} t^{s-1} e^{-t} \, dt$ is the upper incomplete gamma function.

The result is independent of $t_0$, and the choice $t_0 = 1$ typically gives optimal convergence.
\end{theorem}

\subsection{Application to $\Delta$: Explicit Formulas}

For $\Delta(z) = \sum_{n=1}^{\infty} \tau(n) q^n$ at weight $k=12$, choosing $t_0 = 1$:

\begin{proposition}[BFK Expression for $L(\Delta, s)$]
\[
L^*_\Delta(s) = \sum_{n=1}^{\infty} \tau(n) \frac{\Gamma(s, 2\pi n)}{(2\pi n)^s} + \sum_{n=1}^{\infty} \tau(n) \frac{\Gamma(12-s, 2\pi n)}{(2\pi n)^{12-s}}.
\]

Since $L(\Delta, s) = \frac{\Gamma(s)}{(2\pi)^s} L^*_\Delta(s)$, we have:
\[
L(\Delta, s) = \frac{\Gamma(s)}{(2\pi)^s} \left[ \sum_{n=1}^{\infty} \tau(n) \frac{\Gamma(s, 2\pi n)}{(2\pi n)^s} + \sum_{n=1}^{\infty} \tau(n) \frac{\Gamma(12-s, 2\pi n)}{(2\pi n)^{12-s}} \right].
\]
\end{proposition}

\subsection{Computational Test: Verifying Rationality Patterns}

The BFK formula allows direct numerical verification of the lockstep rationality. For instance, at $s = 2$ (even):

\begin{equation}
\frac{L(\Delta, 2)}{(2\pi i)^2 \cdot 1! \cdot \omega_-} \stackrel{?}{=} 768
\end{equation}

We compute:
\begin{align*}
L(\Delta, 2) &= \frac{\Gamma(2)}{(2\pi)^2} \left[ \sum_{n=1}^{\infty} \tau(n) \frac{\Gamma(2, 2\pi n)}{(2\pi n)^2} + \sum_{n=1}^{\infty} \tau(n) \frac{\Gamma(10, 2\pi n)}{(2\pi n)^{10}} \right] \\
&= \frac{1}{(2\pi)^2} \left[ \sum_{n=1}^{\infty} \tau(n) \frac{(1 + 2\pi n) e^{-2\pi n}}{(2\pi n)^2} + \sum_{n=1}^{\infty} \tau(n) \frac{\Gamma(10, 2\pi n)}{(2\pi n)^{10}} \right],
\end{align*}

where we used the recurrence $\Gamma(2, x) = e^{-x}(1 + x)$.

The key computational advantage: the terms decay like $e^{-2\pi n}$, so 10-20 terms typically give 30+ digits of accuracy.

\subsection{Testing Code Structure}

A computational test would involve:

\begin{enumerate}
\item Compute $\omega_-$ from known values of $\zeta(11)$ and $\langle \Delta, \Delta \rangle$
\item Evaluate the BFK sum for $L(\Delta, 2)$ to high precision
\item Check whether $\frac{L(\Delta, 2)}{(2\pi i)^2 \cdot \omega_-} = 768$ to numerical precision
\item Repeat for all critical values $s = 1, \ldots, 11$
\end{enumerate}

\section{Weakly Holomorphic Forms: The BFK Framework}

\subsection{Regularized L-Functions}

For weakly holomorphic forms with poles at cusps, BFK introduce a regularization:

\begin{definition}[BFK Regularized L-Function]
For $f \in M_k^!$ (weakly holomorphic of weight $k$), define:
\[
L^{\text{reg}}_f(s) := \frac{(2\pi)^s}{\Gamma(s)} \mathrm{R.} \int_0^{i\infty} f(iy) y^{s-1} \, dy,
\]
where $\mathrm{R.}$ denotes regularization that removes divergences from the pole.
\end{definition}

\subsection{Period Polynomials for Weakly Holomorphic Forms}

\begin{theorem}[BFK, adapted from Zagier]
For $f \in M_k^!$, there exists an Eichler integral $E_f(z)$ satisfying:
\[
E_f|_{2-k}\gamma - E_f(z) = r_f(z; \gamma)
\]
for $\gamma \in \mathrm{PSL}_2(\mathbb{Z})$, where $r_f(z; \gamma)$ is a period polynomial encoding the regularized L-values.
\end{theorem}

\subsection{Application to $\Delta'$: The Computational Challenge}

The weakly holomorphic form $\Delta' = f_{12,1}$ has Fourier expansion:
\[
\Delta'(z) = q^{-1} + 0 + 0 + c_2 q^2 + c_3 q^3 + \cdots
\]

\begin{proposition}[Expected Period Structure]
The period polynomial $r_{\Delta'}(X)$ should decompose as:
\[
r_{\Delta'}(X) = \tilde{\omega}_+ (\alpha_0 P_0(X) + \alpha_1 P_1(X)) + \tilde{\omega}_- \cdot \alpha_2 P_2(X),
\]
where $\tilde{\omega}_\pm$ are quasi-periods related to regularized L-values.
\end{proposition}

\section{BFK Computational Framework for $\Delta'$}

\subsection{Regularized L-Function Expression}

For $\Delta'(z) = q^{-1} + \sum_{n=2}^{\infty} c_n q^n$, the BFK regularization gives:

\begin{theorem}[BFK Formula for $\Delta'$]
For any $t_0 > 0$:
\[
L^{\text{reg}}_{\Delta'}(s) = \sum_{n=2}^{\infty} c_n \frac{\Gamma(s, 2\pi n t_0)}{(2\pi n)^s} + \sum_{n=2}^{\infty} c_n \frac{\Gamma(12-s, 2\pi n/t_0)}{(2\pi n)^{12-s}} + \text{(pole contribution)},
\]
where the pole contribution from the $q^{-1}$ term must be regularized.
\end{theorem}

\subsection{The Pole Regularization Issue}

The $q^{-1}$ term contributes a divergent piece to the naive integral:
\[
\int_0^{i\infty} q^{-1} y^{s-1} \, dy = \int_0^{\infty} e^{-2\pi y} y^{s-1} \, dy = \frac{\Gamma(s)}{(2\pi)^s}.
\]

This is regular for $\mathrm{Re}(s) > 0$, but the BFK regularization handles this systematically by:
\begin{enumerate}
\item Splitting the integral at a cutoff $t_0$
\item Subtracting the pole contribution explicitly
\item Taking the limit as the regularization parameter vanishes
\end{enumerate}

\begin{remark}
The detailed regularization procedure for $\Delta'$ requires computing:
\[
\mathrm{R.} \int_0^{i\infty} \Delta'(iy) y^{s-1} \, dy = \lim_{\epsilon \to 0} \left[ \int_\epsilon^{\infty} \Delta'(iy) y^{s-1} \, dy - \text{(pole term)} \right].
\]

The BFK paper shows this limit exists and yields well-defined regularized L-values.
\end{remark}

\subsection{Computational Strategy for $\Delta'$}

\subsubsection{Step 1: Computing Fourier Coefficients}

The coefficients $c_n$ of $\Delta'$ can be computed using:

\begin{enumerate}
\item \textbf{Differential equation method}: $\Delta'$ satisfies modular differential equations
\item \textbf{Duke-Jenkins basis relations}: Explicit formulas from the basis construction
\item \textbf{Recurrence relations}: From the modular structure
\end{enumerate}

\begin{proposition}[Fourier Coefficients of $\Delta'$]
The coefficients grow like $c_n \sim e^{4\pi\sqrt{n}}$ (from the circle method), so the BFK sums still converge rapidly despite this growth---the incomplete gamma functions decay like $e^{-2\pi n}$, which dominates.
\end{proposition}

\subsubsection{Step 2: Evaluating Regularized L-Values}

For $s$ in the critical strip ($1 \leq s \leq 11$ for weight 12), we compute:

\begin{equation}
L^{\text{reg}}_{\Delta'}(s) = \frac{(2\pi)^s}{\Gamma(s)} \left[ \sum_{n=2}^{N_{\text{cutoff}}} c_n \frac{\Gamma(s, 2\pi n)}{(2\pi n)^s} + \sum_{n=2}^{N_{\text{cutoff}}} c_n \frac{\Gamma(12-s, 2\pi n)}{(2\pi n)^{12-s}} \right],
\end{equation}

where $N_{\text{cutoff}} \approx 20$-$50$ should suffice for high precision.

\subsubsection{Step 3: Testing Rationality of Quasi-Periods}

\begin{question}[Key Computational Test]
Do the regularized L-values of $\Delta'$ exhibit lockstep rationality analogous to $\Delta$? That is, are the ratios
\[
\frac{L^{\text{reg}}_{\Delta'}(n+1)}{(2\pi i)^{n+1} \cdot n! \cdot \tilde{\omega}_\pm}
\]
rational numbers with controlled denominators?
\end{question}

The computational procedure:
\begin{enumerate}
\item Compute $L^{\text{reg}}_{\Delta'}(2), L^{\text{reg}}_{\Delta'}(4), \ldots, L^{\text{reg}}_{\Delta'}(10)$ (even values)
\item Check if they're all rational multiples of a single period $\tilde{\omega}_-$
\item Compute $L^{\text{reg}}_{\Delta'}(1), L^{\text{reg}}_{\Delta'}(3), \ldots, L^{\text{reg}}_{\Delta'}(11)$ (odd values)
\item Check if they're all rational multiples of a single period $\tilde{\omega}_+$
\item Verify whether the rational coefficients have controlled denominators
\end{enumerate}

\subsection{Anticipated Results and Obstructions}

\subsubsection{What We Expect}

Based on the analogy with $\Delta$, we anticipate:

\begin{enumerate}
\item \textbf{Two-period structure}: Regularized L-values should split into two families, each a rational span of a single quasi-period
\item \textbf{Controlled denominators}: The rational coefficients should involve small primes (2, 3, possibly 691)
\item \textbf{Symmetry with $\Delta$}: The quasi-periods $\tilde{\omega}_\pm$ may have simple relationships to the classical periods $\omega_\pm$
\end{enumerate}

\subsubsection{Potential Complications}

Several factors could complicate the picture:

\begin{enumerate}
\item \textbf{Regularization ambiguity}: Different regularization schemes might yield quasi-periods differing by transcendental factors
\item \textbf{Pole contribution}: The $q^{-1}$ term may introduce additional transcendental constants
\item \textbf{Missing coefficients}: The gap at $q^0$ and $q^1$ breaks the pattern from $\Delta$ and might affect period relations
\item \textbf{Numerical precision}: High-precision arithmetic required due to cancellations in the regularization
\end{enumerate}

\section{Duke-Jenkins Duality and the Hecke Partner}

\subsection{Weakly Holomorphic Modular Form Basis}

\begin{definition}
The Duke-Jenkins basis element $f_{12,m}$ is the unique weakly holomorphic form with:
\begin{itemize}
\item Principal part $q^{-m}$ at the cusp
\item Appropriate growth conditions
\end{itemize}
\end{definition}

At weight 12:
\begin{itemize}
\item $f_{12,-1} = \Delta$ (the classical cusp form, with $q^{-(-1)} = q$ as leading term)
\item $f_{12,1} = \Delta'$ (the Hecke partner, with principal part $q^{-1}$)
\end{itemize}

\subsection{Duality of Coefficients}

\begin{theorem}[Duke-Jenkins]
The Fourier coefficients satisfy:
\[
a_{12}(m,n) = -a_{-10}(n,m),
\]
relating weight 12 weakly holomorphic forms to weight $-10$ forms.
\end{theorem}

This duality suggests deep relationships between periods at weight 12 and periods at weight $-10$.

\section{Computational Investigation of $\Delta'$: Where Classical Theory Breaks Down}

\subsection{Executive Summary of Findings}

We successfully implemented the BFK framework for $\Delta'$ and computed its L-values and period polynomial coefficients. The numerics converge correctly and produce well-defined, structured results. However, these results reveal a fundamental breakdown of the classical period polynomial framework when extended naively to weakly holomorphic forms.

\subsubsection{What Works}
\begin{enumerate}
\item \textbf{BFK formula converges}: The incomplete gamma sum with $\sum_{n \neq 0}$ (including the principal part at $n=-1$) produces stable L-values
\item \textbf{L-values well-defined}: All critical values $L(\Delta', s)$ for $s = 1, \ldots, 11$ are numerically stable
\item \textbf{Zagier's formula applies}: The period polynomial formula works mechanically to extract coefficients
\item \textbf{Structural proportionality}: Coefficients separate cleanly into odd/even sectors with consistent internal ratios
\end{enumerate}

\subsubsection{What Fails}
\begin{enumerate}
\item \textbf{Template polynomial mismatch}: Coefficient ratios are \emph{fundamentally different} from $P_0, P_1, P_2$
\item \textbf{No classical decomposition}: The expected structure $r_{\Delta'}(X) = \tilde{\omega}_+(\alpha_0 P_0 + \alpha_1 P_1) + \tilde{\omega}_- \alpha_2 P_2$ does not hold
\item \textbf{Rationality breaks}: Ratios show $\sim 10^{-6}$ deviations from nearby rationals (vs.\ $< 10^{-27}$ for $\Delta$)
\end{enumerate}

\subsection{Construction of $\Delta'$}

Using the Duke-Jenkins formula, we constructed:
\[
\Delta'(z) = \Delta(z) \cdot (j^2(z) - 1464 j(z) + 142236)
\]

This yields the q-expansion:
\[
\Delta'(z) = q^{-1} + 0 \cdot q^0 + 0 \cdot q^1 + 47709536 q^2 + 39862705122 q^3 + \cdots
\]

The gap structure (no $q^0, q^1$ terms) is verified numerically.

\subsection{BFK L-Values for $\Delta'$: Successful Computation}

The BFK formula with correct regularization (using naked $n$'s, not $|n|$):
\[
L^*_{\Delta'}(s) = \sum_{n \neq 0} c_n \frac{\Gamma(s, 2\pi n)}{(2\pi n)^s} + i^{12} \sum_{n \neq 0} c_n \frac{\Gamma(12-s, 2\pi n)}{(2\pi n)^{12-s}}
\]

For $n = -1$ (the pole), we use $2\pi(-1) = -2\pi$ directly, introducing phase factors.

\begin{table}[h]
\centering
\begin{tabular}{cll}
\hline
$s$ & $L(\Delta, s)$ & $L(\Delta', s)$ \\
\hline
1 & $3.744 \times 10^{-2}$ & $-72.695$ \\
2 & $1.464 \times 10^{-1}$ & $-269.337$ \\
3 & $3.152 \times 10^{-1}$ & $-581.852$ \\
$\vdots$ & $\vdots$ & $\vdots$ \\
11 & $9.894 \times 10^{-1}$ & $-1921.050$ \\
\hline
\end{tabular}
\caption{L-values for $\Delta$ vs $\Delta'$ (with corrected BFK formula). Note: $\Delta'$ values are negative and 2-3 orders of magnitude larger.}
\end{table}

\subsection{Period Polynomial Coefficients: Side-by-Side Comparison}

Applying Zagier's formula yields period polynomial coefficients. The critical finding: \textbf{coefficient ratios for $\Delta'$ do not match the classical templates}.

\subsubsection{Odd Part: $P_2$ Template Mismatch}

\begin{table}[h]
\centering
\begin{tabular}{cccccc}
\hline
Power & \multicolumn{2}{c}{$\Delta$ (holomorphic)} & \multicolumn{2}{c}{$\Delta'$ (weakly holo.)} & Match? \\
\cline{2-5}
& Coeff & Ratio & Coeff & Ratio & \\
\hline
$X^9$ & $0.0371$ & $\mathbf{1.000}$ & $985.1$ & $\mathbf{1.000}$ & -- \\
$X^7$ & $-0.2317$ & $-6.250$ & $-10864.8$ & $-11.029$ & \textbf{NO} \\
$X^5$ & $0.3893$ & $10.500$ & $22208.0$ & $22.544$ & \textbf{NO} \\
$X^3$ & $-0.2317$ & $-6.250$ & $-10864.8$ & $-11.029$ & \textbf{NO} \\
$X^1$ & $0.0371$ & $\mathbf{1.000}$ & $985.1$ & $\mathbf{1.000}$ & -- \\
\hline
\end{tabular}
\caption{Odd power comparison. $\Delta$ matches $P_2$ template $(1, -25/4, 21/2, -25/4, 1)$ to machine precision. $\Delta'$ has \emph{different polynomial} $(1, -11.029, 22.544, -11.029, 1)$.}
\end{table}

\subsubsection{Even Part: $P_0 + P_1$ Template Mismatch}

\begin{table}[h]
\centering
\begin{tabular}{cccccc}
\hline
Power & \multicolumn{2}{c}{$\Delta$ (holomorphic)} & \multicolumn{2}{c}{$\Delta'$ (weakly holo.)} & Match? \\
\cline{2-5}
& Coeff & Ratio & Coeff & Ratio & \\
\hline
$X^{10}$ & $0.0060\,i$ & $\mathbf{1.000}$ & $105.3\,i$ & $\mathbf{1.000}$ & -- \\
$X^8$ & $-0.1144\,i$ & $-19.194$ & $-4217.5\,i$ & $-40.051$ & \textbf{NO} \\
$X^6$ & $0.3431\,i$ & $57.583$ & $18631.1\,i$ & $176.931$ & \textbf{NO} \\
$X^4$ & $-0.3431\,i$ & $-57.583$ & $-18631.1\,i$ & $-176.931$ & \textbf{NO} \\
$X^2$ & $0.1144\,i$ & $19.194$ & $4217.5\,i$ & $40.051$ & \textbf{NO} \\
$X^0$ & $-0.0060\,i$ & $\mathbf{-1.000}$ & $-105.3\,i$ & $\mathbf{-1.000}$ & -- \\
\hline
\end{tabular}
\caption{Even power comparison. $\Delta$ matches $P_0 + P_1$ with 691 appearing naturally. $\Delta'$ shows completely different ratios.}
\end{table}

\subsection{Quantitative Analysis: Not Numerical Noise}

The differences are \emph{not} small perturbations or numerical errors:

\begin{table}[h]
\centering
\begin{tabular}{lcccc}
\hline
Sector & Coefficient & $\Delta$ ratio & $\Delta'$ ratio & Factor diff. \\
\hline
Odd & $c_7/c_9$ & $-6.25$ & $-11.029$ & $\mathbf{1.76\times}$ \\
Odd & $c_5/c_9$ & $10.5$ & $22.544$ & $\mathbf{2.15\times}$ \\
Even & $c_8/c_{10}$ & $-19.194$ & $-40.051$ & $\mathbf{2.09\times}$ \\
Even & $c_6/c_{10}$ & $57.583$ & $176.931$ & $\mathbf{3.07\times}$ \\
\hline
\end{tabular}
\caption{Factors of 1.5--3.0 indicate fundamentally different template polynomials, not numerical noise.}
\end{table}

\subsection{Rationality Patterns: Breakdown of Classical Structure}

\begin{table}[h]
\centering
\begin{tabular}{lccc}
\hline
Form & Sector & Rationality Error & Interpretation \\
\hline
$\Delta$ & Odd & $< 10^{-27}$ & Machine precision rational \\
$\Delta$ & Even & $< 10^{-27}$ & Machine precision rational \\
\hline
$\Delta'$ & Odd & $\sim 10^{-6}$ & Forced fit artifacts \\
$\Delta'$ & Even & $\sim 10^{-6}$ & Forced fit artifacts \\
\hline
\end{tabular}
\caption{For $\Delta$, ratios are rational to machine precision because they match known templates $P_0, P_1, P_2$. For $\Delta'$, the $\sim 10^{-6}$ ``errors'' are artifacts of forcing floats into nearby rationals when the true structure is \emph{different}.}
\end{table}

\subsection{Mathematical Debugging: Where the Framework Breaks}

\subsubsection{Classical Framework Assumptions}

Zagier's period polynomial theory for holomorphic cusp forms relies on:
\begin{enumerate}
\item Eichler-Shimura isomorphism: $H^1_{\text{par}}(X_0(N), \mathbb{C}) \cong S_k \oplus \overline{S_k}$
\item Period integrals $\int_0^{i\infty} f(\tau) \tau^n d\tau$ encode cohomology classes
\item Functional equation: $r_f(X) + (-1)^k r_f(-X) = 0$ (for cusp forms)
\item Template basis: $P_j$ span the space of period polynomials modulo functional equations
\end{enumerate}

\subsubsection{What Changes for $\Delta'$?}

\begin{enumerate}
\item \textbf{Period integrals diverge}: $\int_0^{i\infty} \Delta'(\tau) \tau^n d\tau$ has pole contribution
\item \textbf{Cohomology unclear}: What replaces $H^1_{\text{par}}$ for forms with growth at cusps?
\item \textbf{Eichler-Shimura}: Does the isomorphism extend? With what modifications?
\item \textbf{Template basis}: Are there \emph{different} template polynomials $\tilde{P}_j$ for weakly holomorphic forms?
\end{enumerate}

\subsubsection{Potential Frameworks}

Several cohomological setups might explain the $\Delta'$ structure:

\begin{enumerate}
\item \textbf{Parabolic cohomology with growth conditions}: $H^1_{\text{par}}(X, *)$ allowing controlled growth
\item \textbf{Relative cohomology}: $H^1(X, \{\text{cusps}\}, \mathbb{C})$ accounting for boundary behavior
\item \textbf{Mixed Hodge structures}: For non-compact modular curves
\item \textbf{Regularized period theory}: Systematic analytic continuation of divergent integrals
\item \textbf{Higher-order modular forms}: Framework of Bruggeman-Choie-Diamantis
\end{enumerate}

\subsection{Key Research Questions}

\begin{question}[Template Polynomial Basis]
Does $\Delta'$ admit a period polynomial decomposition using \emph{different} templates $\tilde{P}_0, \tilde{P}_1, \tilde{P}_2$?
If so, what is the relationship between $\{P_j\}$ and $\{\tilde{P}_j\}$?
\end{question}

The numerical evidence suggests:
\begin{itemize}
\item Odd part follows $(1, -11.029, 22.544, -11.029, 1)$ $\to$ some polynomial $\tilde{P}_2 \neq P_2$
\item Even part follows $(1, -40.051, 176.931, \ldots)$ $\to$ combination $\tilde{P}_0 + \tilde{P}_1 \neq P_0 + P_1$
\end{itemize}

\begin{question}[BFK to Periods Map]
Is Zagier's formula the correct bridge from BFK L-values to periods for weakly holomorphic forms?
Or does the pole regularization require a \emph{modified formula}?
\end{question}

The BFK L-values are well-defined, but their relationship to cohomological periods may differ from the holomorphic case.

\begin{question}[Quasi-Periods]
What are the quasi-periods $\tilde{\omega}_\pm$ for $\Delta'$, and how do they relate to the classical periods $\omega_\pm$ of $\Delta$?
\end{question}

\subsection{Computational Artifacts}

A minimal self-contained demonstration is provided in \texttt{minimal\_delta\_comparison.sage}, which:
\begin{itemize}
\item Implements BFK with corrected formula (naked $n$'s, no absolute values)
\item Computes period polynomials for both $\Delta$ and $\Delta'$
\item Shows $\Delta$ matches classical templates to $< 10^{-27}$ precision
\item Shows $\Delta'$ has fundamentally different polynomial structure (factors of 1.5--3.0 away)
\end{itemize}

\subsection{Conclusion}

\textbf{The numerics work.} The BFK framework successfully computes L-values for $\Delta'$, and Zagier's formula mechanically produces period polynomial coefficients.

\textbf{The mathematics breaks down.} The classical period polynomial theory---designed for holomorphic cusp forms via Eichler-Shimura cohomology---does not extend naively to weakly holomorphic forms. The coefficient structure suggests $\Delta'$ has its own template polynomial basis, distinct from $\{P_0, P_1, P_2\}$.

\textbf{We need better theory.} Progress requires understanding:
\begin{itemize}
\item The correct cohomological framework for weakly holomorphic forms
\item How pole regularization affects period structures
\item Whether different template polynomials exist and how to characterize them
\item The relationship between BFK regularization and cohomological regularization
\end{itemize}

\textbf{The real question is not} ``Why don't the numerics give rational ratios?'' but rather ``What cohomological framework explains what the numerics are telling us?''

\section{Open Questions and Computational Research Program}

\subsection{Period Structure of $\Delta'$}

\begin{question}
What are the explicit values (or relationships to $\omega_\pm$) of the coefficients $\alpha_0, \alpha_1, \alpha_2$ in the period polynomial decomposition of $\Delta'$?
\end{question}

Expected structure:
\[
r_{\Delta'}(X) = \tilde{\omega}_+ (\alpha_0 P_0(X) + \alpha_1 P_1(X)) + \tilde{\omega}_- \cdot \alpha_2 P_2(X)
\]

\subsection{Rationality of Quasi-Periods}

\begin{question}
Do the regularized L-values of $\Delta'$ exhibit the same lockstep rationality as those of $\Delta$? That is, are the ratios
\[
\frac{L^{\text{reg}}_{\Delta'}(n+1)}{(2\pi i)^{n+1} \cdot n! \cdot \tilde{\omega}_\pm}
\]
rational numbers with controlled denominators?
\end{question}

\subsection{Computational Research Program}

To address these questions systematically:

\subsubsection{Phase 1: Fourier Coefficients of $\Delta'$}

\begin{enumerate}
\item Implement computation of $c_n$ for $n = 2, 3, \ldots, 100$ using Duke-Jenkins formulas or differential equations
\item Verify growth rate matches theoretical prediction $c_n \sim e^{4\pi\sqrt{n}}$
\item Check for any additional arithmetic patterns in the coefficients
\end{enumerate}

\subsubsection{Phase 2: Regularized L-Values}

\begin{enumerate}
\item Implement BFK incomplete gamma sum for $L^{\text{reg}}_{\Delta'}(s)$
\item Compute all critical values $s = 1, 2, \ldots, 11$ to at least 50 digits precision
\item Test convergence rate: verify that 20-50 terms suffice for target precision
\end{enumerate}

\subsubsection{Phase 3: Rationality Testing}

\begin{enumerate}
\item For even $s$: compute ratios $L^{\text{reg}}_{\Delta'}(s) / L^{\text{reg}}_{\Delta'}(2)$
\begin{itemize}
\item Use PSLQ or LLL to detect rational relationships
\item Check if denominators involve only small primes
\end{itemize}
\item For odd $s$: compute ratios $L^{\text{reg}}_{\Delta'}(s) / L^{\text{reg}}_{\Delta'}(1)$
\begin{itemize}
\item Test for rationality
\item Compare denominator structure to $\Delta$ case
\end{itemize}
\item Test for relationships between $\tilde{\omega}_+$ and $\tilde{\omega}_-$
\item Investigate connections to classical periods $\omega_\pm$ of $\Delta$
\end{enumerate}

\subsubsection{Phase 4: Period Polynomial Reconstruction}

\begin{enumerate}
\item Extract coefficients $\alpha_0, \alpha_1, \alpha_2$ from the L-value data
\item Verify period polynomial relations hold
\item Compare to predictions from Duke-Jenkins duality
\end{enumerate}

\subsection{Relationship to Higher Weights}

\begin{question}
For weights $k > 12$ where $\dim S_k > 1$, how do multiple Hecke eigenforms contribute to the period polynomial basis? Do rationality patterns persist, but with coefficients in Hecke fields rather than $\mathbb{Q}$?
\end{question}

\subsection{Anticipated Challenges}

\subsubsection{Numerical Issues}

\begin{enumerate}
\item \textbf{Precision requirements}: Regularization may involve significant cancellation, requiring 100+ digit arithmetic
\item \textbf{Coefficient growth}: $c_n \sim e^{4\pi\sqrt{n}}$ is rapid; need careful error analysis
\item \textbf{Incomplete gamma evaluation}: Requires robust implementation for large arguments
\end{enumerate}

\subsubsection{Theoretical Issues}

\begin{enumerate}
\item \textbf{Regularization scheme dependence}: Different regularizations might give different quasi-periods
\item \textbf{Normalization conventions}: Choice of normalization for $\Delta'$ affects all computations
\item \textbf{Interpretation}: Even if rationality holds numerically, understanding \emph{why} requires deeper theory
\end{enumerate}

\section{Connection to Harmonic Maass Forms (Optional)}

\begin{remark}[Harmonic Maass Forms]
The BFK paper discusses connections to harmonic weak Maass forms: forms $F$ of weight $2-k$ whose derivatives (via $\xi$ operator) give cusp forms. These have ``mock periods'' $c_F^\pm(\pm 1)$ that satisfy limiting relationships with classical L-values.

However, harmonic Maass forms are \textbf{not required} to understand the period structure of $\Delta$ and $\Delta'$. The quasi-periods of $\Delta'$ arise directly from its regularized L-values through the period polynomial machinery, independent of any harmonic Maass construction.

We mention this only for completeness, as the BFK paper discusses it, but it is tangential to the core computational program outlined here.
\end{remark}

\section{Conclusion}

The weight 12 case, with its unique cusp form $\Delta$ and canonical weakly holomorphic partner $\Delta'$, provides an ideal testing ground for understanding:

\begin{itemize}
\item The rigid rationality structure of critical L-values
\item Extensions of period polynomial theory to weakly holomorphic forms with poles at cusps  
\item Connections between classical periods and quasi-periods of Hecke partners
\item Duality structures relating positive and negative weight forms
\end{itemize}

The BFK incomplete gamma framework provides the computational machinery to test these conjectures numerically with high precision. The key open question is whether the quasi-periods of $\Delta'$ exhibit the same lockstep rational structure as the classical periods of $\Delta$.

\subsection{Summary of Computational Approach}

The BFK formulas transform what would be slowly convergent Dirichlet series into exponentially convergent sums:

\begin{center}
\begin{tabular}{l|l}
\textbf{Classical} & \textbf{BFK Incomplete Gamma} \\
\hline
$L(f,s) = \sum \frac{a_n}{n^s}$ & $L^*_f(s) = \sum a_n \frac{\Gamma(s, 2\pi n t_0)}{(2\pi n)^s} + \cdots$ \\
Convergence: $\sim 1/n^s$ & Convergence: $\sim e^{-2\pi n}$ \\
$\Rightarrow$ slow & $\Rightarrow$ exponentially fast \\
\end{tabular}
\end{center}

This makes high-precision numerical verification of period rationality feasible, potentially leading to:
\begin{enumerate}
\item Numerical discovery of the exact rational coefficients
\item Conjectures about the structure for higher weights
\item Insight into the regularization mechanism for weakly holomorphic forms
\end{enumerate}

\begin{thebibliography}{9}

\bibitem{zagier}
Don Zagier, \textit{Periods of modular forms and Jacobi theta functions}, Inventiones mathematicae \textbf{104} (1991), 449--465.

\bibitem{bfk}
Kathrin Bringmann, Karl-Heinz Fricke, and Zachary A. Kent, \textit{Special L-values and periods of weakly holomorphic modular forms}, Ramanujan Journal \textbf{24} (2011), 119--139.

\bibitem{duke_jenkins}
William Duke and Paul Jenkins, \textit{On the zeros and coefficients of certain weakly holomorphic modular forms}, Pure and Applied Mathematics Quarterly \textbf{4} (2008), 1327--1340.

\bibitem{eichler}
Martin Eichler, \textit{Eine Verallgemeinerung der Abelschen Integrale}, Mathematische Zeitschrift \textbf{67} (1957), 267--298.

\bibitem{shimura}
Goro Shimura, \textit{Sur les intégrales attachées aux formes automorphes}, Journal of the Mathematical Society of Japan \textbf{11} (1959), 291--311.

\bibitem{brown}
Francis Brown, \textit{Multiple modular values and the relative completion of the fundamental group of $\mathcal{M}_{1,1}$}, arXiv:1407.5167 (2017).

\bibitem{petersen}
Kathrin Petersen, \textit{Eichler cohomology and zeros of polynomials associated to derivatives of L-functions}, arXiv:1710.07912 (2017).

\bibitem{bruggeman}
Roelof Bruggeman, YoungJu Choie, and Nikolaos Diamantis, \textit{Holomorphic automorphic forms and cohomology}, Memoirs of the American Mathematical Society \textbf{253} (2018), no. 1212.

\end{thebibliography}

\end{document}
